\documentclass{article}
\usepackage[utf8]{inputenc}
\usepackage{amsmath,amssymb}
\usepackage[most]{tcolorbox}
\usepackage{geometry}
\geometry{margin=2.5cm}

\newcommand{\step}[1]{\par\noindent\hangindent=3.5em\hangafter=1%
  \makebox[3.5em][l]{\textbf{#1.}}}
\newcommand{\steptag}[1]{\hfill{\small\textsf{[#1]}}}

\newtcolorbox{proofbox}[1]{
  colback=gray!5, colframe=gray!60,
  fonttitle=\bfseries, title={#1},
  breakable, enhanced,
  left=4pt, right=4pt, top=4pt, bottom=4pt,
  before skip=10pt, after skip=10pt
}

\begin{document}

\begin{proofbox}{Uniform global tangent/Maurer-Cartan control: there are u...}
\step{1} Uniform global tangent/Maurer-Cartan control: there are universal C\_ch >= 1, kappa\_ch in (0, 1/2] such that for every finite-dimensional exact-unit epsilon\_r-C*-algebra and every delta > 0 with C\_ch*(epsilon\_r + delta) <= kappa\_ch, the graph maps supplied by lem-stage1-unitary-graph-control satisfy: every tangent space T\_U calU is the image of L\_U(I + Dg\_U(0)): icalH -> calX, and omega\_U(Z) = (L\_U\textasciicircum{}\{-1\} Z)\textasciicircum{}par : T\_U calU -> icalH is a global C\textasciicircum{}1 bundle trivialization with distortion at most 1 + C\_ch*epsilon\_r, satisfying omega\_\{cU\}(cZ) = omega\_U(Z) and omega\_U(iU) = iJ for every c in U(1). \steptag{validated} \\[6pt]
\step{1.1} Witness and guard selection: let C0 >= 1 and kappa0 in (0,1/2] be the universal witnesses of lem-stage1-unitary-graph-control. There is a universal K >= 1 such that, whenever 0 <= epsilon\_r <= 1/4 and U is in calU, one has ||U|| <= (1-epsilon\_r)\textasciicircum{}(-1/2), ||L\_U|| <= 1+K*epsilon\_r, and ||L\_U\textasciicircum{}\{-1\}|| <= 1+K*epsilon\_r; choose C\_ch >= max(C0,K,4) large enough to absorb all fixed products of these bounds and the graph bound, and choose 0 < kappa\_ch <= min(kappa0,1/4). Then C\_ch*(epsilon\_r+delta) <= kappa\_ch implies the upstream graph guard and epsilon\_r <= 1/4. \steptag{validated} \\[6pt]
\step{1.1.1} Approximate left multiplication estimate from def-epsilon-cstar-algebra and def-approximate-unitary-space: for U in calU, U\textasciicircum{}dagger bold-dot U=J and ||J||=1. The lower C*-inequality gives ||U||=||U\textasciicircum{}dagger|| <= (1-epsilon\_r)\textasciicircum{}(-1/2), and the product inequality gives ||L\_U||,||L\_\{U\textasciicircum{}dagger\}|| <= (1+epsilon\_r)||U||. The associator axiom and exact unit give ||L\_\{U\textasciicircum{}dagger\}L\_U-I|| <= epsilon\_r||U||\textasciicircum{}2 <= epsilon\_r/(1-epsilon\_r). For epsilon\_r <= 1/4 this is <1, so the Neumann lemma and finite dimensionality make L\_U invertible and give ||L\_U\textasciicircum{}\{-1\}|| <= (1-epsilon\_r/(1-epsilon\_r))\textasciicircum{}(-1)*(1+epsilon\_r)*(1-epsilon\_r)\textasciicircum{}(-1/2). These displayed scalar functions are at most 1+K*epsilon\_r on [0,1/4] for a universal K, and the same is true of ||L\_U||. \steptag{validated} \\[4pt]
\step{1.1.2} Constant absorption and application of lem-stage1-unitary-graph-control: take its universal witnesses C0,kappa0. Choose a universal C\_ch >= max(C0,K,4) large enough that every fixed product used in the tangent distortion is bounded by 1+C\_ch*epsilon\_r for 0 <= epsilon\_r <= 1/4, and set kappa\_ch=min(kappa0,1/4). If C\_ch*(epsilon\_r+delta)<=kappa\_ch with delta>0, then C0*(epsilon\_r+delta)<=kappa0 and epsilon\_r<1/4, so the external graph theorem applies; enlarging C\_ch once more does not spoil this implication. \steptag{validated} \\[4pt]
\step{1.2} Centered graph differentiation: for every admissible algebra and U in calU, the graph map of lem-stage1-unitary-graph-control centered at U satisfies g\_U(0)=0 and ||Dg\_U(0)|| <= C0*epsilon\_r, and its graph-chart parametrization has derivative L\_U(I+Dg\_U(0)) at 0; hence T\_U calU is exactly the image of that map. \steptag{validated} \\[6pt]
\step{1.2.1} Basepoint and sharp derivative bound using lem-stage1-unitary-graph-control: exact unitality and U\textasciicircum{}dagger bold-dot U=J give f\_U(0)=0, so uniqueness gives g\_U(0)=0. If rho>0 also satisfies C0*(epsilon\_r+rho)<=kappa0, the external theorem gives a graph germ g\_U\textasciicircum{}\{rho\} with ||Dg\_U\textasciicircum{}\{rho\}(0)||<=C0*(epsilon\_r+rho). For two such radii, continuity at 0 and g\_U\textasciicircum{}\{rho\}(0)=0 put both solutions in the smaller Hermitian ball on a common neighborhood; pointwise uniqueness there makes the germs, hence their derivatives at 0, equal. Taking rho down to 0 therefore yields ||Dg\_U(0)||<=C0*epsilon\_r for the graph map at the originally prescribed delta. \steptag{validated} \\[4pt]
\step{1.2.2} Tangent calculation using lem-stage1-unitary-graph-control: the centered chart is Psi\_U(A)=U bold-dot (J+A+g\_U(A)) for A in icalH. It obeys Psi\_U(0)=U, is C\textasciicircum{}1, and differentiating the bilinear product gives D Psi\_U(0)A=L\_U(A+Dg\_U(0)A)=L\_U(I+Dg\_U(0))A. Because the external theorem says these are graph charts covering calU, the image of this chart derivative is exactly T\_U calU. \steptag{validated} \\[4pt]
\step{1.3} Fibrewise inverse, distortion, and regularity: on the tangent image, omega\_U(Z)=(L\_U\textasciicircum{}\{-1\}Z)\textasciicircum{}par is inverse to A maps-to L\_U(A+Dg\_U(0)A). The estimates from the witness step and ||Dg\_U(0)|| <= C0*epsilon\_r give norms of omega\_U and its inverse at most 1+C\_ch*epsilon\_r; the displayed omega is a global C\textasciicircum{}1 bundle trivialization. \steptag{validated} \\[6pt]
\step{1.3.1} Detailed inverse/distortion/regularity argument: write P\textasciicircum{}par(X)=(X-X\textasciicircum{}dagger)/2, whose norm is at most 1 by the isometric involution in def-epsilon-cstar-algebra. For A in icalH, Dg\_U(0)A lies in calH, hence P\textasciicircum{}par(A+Dg\_U(0)A)=A. Therefore omega\_U(L\_U(A+Dg\_U(0)A))=A, and the tangent-image result makes omega\_U a fibrewise isomorphism with that displayed inverse. Its norm is at most ||L\_U\textasciicircum{}\{-1\}||, while the inverse has norm at most ||L\_U||*(1+||Dg\_U(0)||); the witness selection absorbs both into 1+C\_ch*epsilon\_r. Finally U maps-to L\_U is linear, inversion is C\textasciicircum{}1 on the open set of invertible endomorphisms, and P\textasciicircum{}par is fixed linear, so (U,Z) maps-to (U,P\textasciicircum{}par L\_U\textasciicircum{}\{-1\}Z) restricts to a C\textasciicircum{}1 tangent-bundle map. A fibrewise bijective C\textasciicircum{}1 bundle map has C\textasciicircum{}1 inverse in local frames because finite-dimensional matrix inversion is C\textasciicircum{}1; thus omega is the asserted global C\textasciicircum{}1 bundle trivialization with the stated distortion. \steptag{validated} \\[6pt]
\step{1.3.1.1} Dependency-closed inverse/distortion/regularity bridge. By validated nodes 1.1.1 and 1.1.2, L\_U is invertible and both ||L\_U|| and ||L\_U\textasciicircum{}\{-1\}|| are at most 1+K*epsilon\_r, with the final witness C\_ch chosen to absorb the fixed products below. By validated nodes 1.2.1 and 1.2.2, g\_U(0)=0, ||Dg\_U(0)||<=C0*epsilon\_r, Dg\_U(0)(icalH) is contained in calH, and T\_U calU is the image of A maps-to L\_U(A+Dg\_U(0)A). Let P\textasciicircum{}par(X)=(X-X\textasciicircum{}dagger)/2. The isometric involution gives ||P\textasciicircum{}par||<=1, and for A in icalH and H in calH one has P\textasciicircum{}par(A+H)=A. Hence P\textasciicircum{}par L\_U\textasciicircum{}\{-1\} sends L\_U(A+Dg\_U(0)A) to A; the tangent-image equality makes it a fibrewise bijection with inverse A maps-to L\_U(A+Dg\_U(0)A). Its two operator norms are bounded respectively by ||L\_U\textasciicircum{}\{-1\}|| and ||L\_U||*(1+||Dg\_U(0)||), hence by 1+C\_ch*epsilon\_r by the choice in 1.1.2. For regularity, U maps-to L\_U is linear, inversion is C\textasciicircum{}1 on the open set of invertible endomorphisms, and P\textasciicircum{}par is fixed real-linear, so (U,Z) maps-to (U,P\textasciicircum{}par L\_U\textasciicircum{}\{-1\}Z) is C\textasciicircum{}1 before restriction and therefore on TcalU. In local bundle frames it is represented by an everywhere-invertible C\textasciicircum{}1 matrix B(U); B(U)\textasciicircum{}\{-1\}=adj(B(U))/det(B(U)) is C\textasciicircum{}1. The local inverses agree with the displayed fibrewise inverse, so they glue to a global C\textasciicircum{}1 inverse. Thus omega is the asserted global C\textasciicircum{}1 trivialization, using only the four declared validated dependencies rather than the pending parents 1.1 and 1.2. \steptag{validated} \\[4pt]
\step{1.3.1.2} Smooth dependence of the graph derivative from the defining equation: for U in calU set B\_U=D\_\{A\textasciicircum{}perp\}f\_U(0)|\_\{calH\}:calH->calH and C\_U=D\_\{A\textasciicircum{}par\}f\_U(0)|\_\{icalH\}:icalH->calH. The upstream formula f\_U(A)=(1/2)*(((J+A\textasciicircum{}dagger) bold-dot U\textasciicircum{}dagger) bold-dot (U bold-dot (J+A))-J), together with fixed bilinear multiplication and the fixed real-linear involution, shows directly that U maps-to B\_U and U maps-to C\_U are polynomial, hence C\textasciicircum{}1, maps into the finite-dimensional operator spaces. Since U is unitary, f\_U(0)=0 and uniqueness gives g\_U(0)=0. The upstream estimate at A\textasciicircum{}par=0 therefore gives ||B\_U-I\_calH||<=C0*epsilon\_r<1, so B\_U is invertible. Differentiating the C\textasciicircum{}1 identity f\_U(A+g\_U(A))=0 at A=0 in each anti-Hermitian direction gives C\_U+B\_U Dg\_U(0)=0, whence Dg\_U(0)=-B\_U\textasciicircum{}\{-1\}C\_U. Inversion on the open set of invertible finite-dimensional operators is C\textasciicircum{}1, so U maps-to Dg\_U(0) is C\textasciicircum{}1 on calU. Consequently (U,A) maps-to (U,L\_U(A+Dg\_U(0)A)) is a C\textasciicircum{}1 inverse to omega; this proves inverse regularity without inferring it from mere C\textasciicircum{}1 graph-chart regularity. \steptag{validated} \\[4pt]
\step{1.3.1.3} Correction to the challenged regularity inference: do not infer C\textasciicircum{}1 regularity of the inverse from a merely C\textasciicircum{}1 graph frame. Instead use node 1.3.1.2: the differentiated defining equation gives the explicit smooth formula Dg\_U(0)=-B\_U\textasciicircum{}\{-1\}C\_U, with B\_U and C\_U polynomial in U and B\_U invertible by the upstream near-identity estimate. Hence the already identified fibrewise inverse (U,A) maps to (U,L\_U(A+Dg\_U(0)A)) is C\textasciicircum{}1 directly. The counterexample in the challenge is inapplicable precisely because its graph derivative has no analogous smooth implicit-equation formula with invertible normal derivative. \steptag{validated} \\[4pt]
\step{1.4} Equivariance and central tangent: for every c in U(1), scalar multiplication carries T\_U calU to T\_\{cU\} calU and omega\_\{cU\}(cZ)=omega\_U(Z); moreover the scalar orbit through U has velocity iU and omega\_U(iU)=iJ. \steptag{validated} \\[6pt]
\step{1.4.1} Scalar equivariance from def-approximate-unitary-space and bilinearity: U maps-to cU preserves calU and is a linear diffeomorphism with derivative Z maps-to cZ, hence cZ lies in T\_\{cU\}calU. Since L\_\{cU\}=cL\_U, one has L\_\{cU\}\textasciicircum{}\{-1\}(cZ)=L\_U\textasciicircum{}\{-1\}Z, and taking the anti-Hermitian component gives omega\_\{cU\}(cZ)=omega\_U(Z). \steptag{validated} \\[4pt]
\step{1.4.2} Central tangent identity from def-approximate-unitary-space and exact unitality: the curve t maps-to exp(it)U lies in calU and has derivative iU at zero, so iU lies in T\_UcalU. Complex bilinearity and U bold-dot J=U give L\_U(iJ)=iU; therefore L\_U\textasciicircum{}\{-1\}(iU)=iJ, which is anti-Hermitian because J\textasciicircum{}dagger=J, and omega\_U(iU)=iJ. \steptag{validated} \\[4pt]
\end{proofbox}

\end{document}
